\clearpage
\section*{Appendix AR. Normalized seven-ray identities and sharp scalar coefficient floor}
\addcontentsline{toc}{section}{Appendix AR. Normalized seven-ray identities and sharp scalar coefficient floor}

\noindent
This appendix records two exact algebraic statements underlying the normalized
seven-ray tensor representation. Their hypotheses are part of each assertion.

\paragraph{The normalized trigonometric tensor identity.}
For real \(r,\ell,s,c\) define
\[
\begin{aligned}
C_4&=2+16s^2+16\ell^2(r^2s^2-r^4c^2)+16\ell rcs,\\
C_5&=2+16s^2+16\ell^2(r^2s^2-r^4c^2)-16\ell rcs,\\
C_6&=1-16\ell rcs,\qquad C_7=1+16\ell rcs,\\
\Pi&=\frac38-2c^2(1+\ell^2r^4).
\end{aligned}
\]
\textbf{Proposition AR.1.}
If \(s^2+c^2=1\), then
\begin{equation}\label{ar:zz}
\frac{C_4+C_5+C_6+C_7}{16}
=2(1+\ell^2r^2s^2)+\Pi.
\end{equation}
\emph{Proof.} Expanding the left-hand side gives
\[
\frac38+2s^2+2\ell^2r^2s^2-2\ell^2r^4c^2.
\]
Subtracting the right-hand side of \eqref{ar:zz} reduces exactly to
\(2(s^2+c^2-1)\), which vanishes under the stated hypothesis. \(\square\)

\paragraph{A sharp scalar floor.}
Let \(\mu,d_1,\ldots,d_7\in\mathbb R\) with \(d_i\geq\mu\) for
each \(1\leq i\leq7\). Define the normalized diagonal functional
\[
\mathcal L(d)=\sum_{i=1}^{7}d_i-
\frac{d_2+d_3+d_6+d_7}{16}.
\]
\textbf{Proposition AR.2.}
For every real \(\mu\),
\begin{equation}\label{ar:sharp-floor}
\mathcal L(d)\geq\frac{27}{4}\mu.
\end{equation}
The factor \(27/4\) is sharp for the stated coordinate floors.
If \(\mu\geq0\), the weaker estimate
\(\mathcal L(d)\geq (105/16)\mu\) follows.

\emph{Proof.}
The functional has coefficient \(1\) at \(d_1,d_4,d_5\)
and coefficient \(15/16\) at \(d_2,d_3,d_6,d_7\).
All seven coefficients are positive, and therefore
\[
\mathcal L(d)\geq
\left(3+4\frac{15}{16}\right)\mu
=\frac{27}{4}\mu.
\]
Equality holds for \(d_1=\cdots=d_7=\mu\), proving sharpness.
For nonnegative \(\mu\), the remaining implication uses
\(\frac{27}{4}-\frac{105}{16}=\frac{3}{16}\geq0\).
\(\square\)

The two propositions are scalar algebraic statements. They apply to a
physical H8 stress field only once the same normalized parameters and
one-sided seven-ray coefficient floors have been established for that field.
